% !TeX program = tectonic
\documentclass[11pt,a4paper]{article}
\usepackage{fontspec}
\usepackage[english]{babel}
\babelfont{rm}[Language=Default]{Liberation Serif}
\babelfont[Language=Default]{sf}{Carlito}
\babelfont[Language=Default]{tt}{DejaVu Sans Mono}
\usepackage{amsmath,amssymb,amsthm}
\usepackage{geometry}
\geometry{a4paper, top=2.4cm, bottom=2.4cm, left=2.4cm, right=2.4cm}
\usepackage{booktabs,tabularx,enumitem,xcolor,tcolorbox}
\tcbuselibrary{breakable,skins}
\definecolor{linkblue}{HTML}{1F4E79}
\definecolor{statgreen}{HTML}{2F6B3A}
\definecolor{goldacc}{HTML}{8B7E5A}
\usepackage[colorlinks=true,linkcolor=linkblue,urlcolor=linkblue,unicode=true]{hyperref}
\theoremstyle{plain}
\newtheorem{theorem}{Theorem}
\newtheorem{lemma}[theorem]{Lemma}
\newtheorem{proposition}[theorem]{Proposition}
\newtheorem{corollary}[theorem]{Corollary}
\theoremstyle{definition}
\newtheorem{definition}[theorem]{Definition}
\theoremstyle{remark}
\newtheorem{remark}[theorem]{Remark}
\newtcolorbox{statusbox}[1][]{colback=statgreen!5,colframe=statgreen!75,
  title={\textbf{Summary of results:} #1},fonttitle=\small,boxrule=0.7pt,arc=2pt,breakable}
\newtcolorbox{databox}[1][]{colback=goldacc!6,colframe=goldacc!80,
  title={\textbf{Protocol data:} #1},fonttitle=\small,boxrule=0.7pt,arc=2pt,breakable}
\newtcolorbox{verifybox}[1][]{colback=linkblue!5,colframe=linkblue!80,
  title={\textbf{Verification:} #1},fonttitle=\small,boxrule=0.7pt,arc=2pt,breakable}
\widowpenalty=10000
\clubpenalty=10000
\setlength{\parskip}{0.35em}
\newcommand{\CC}{\mathbb{C}}
\newcommand{\RR}{\mathbb{R}}
\newcommand{\ZZ}{\mathbb{Z}}
\newcommand{\QQ}{\mathbb{Q}}
\newcommand{\Ga}{\Gamma}
\newcommand{\Bfun}{\mathrm{B}}
\newcommand{\im}{\operatorname{Im}}
\newcommand{\re}{\operatorname{Re}}
\newcommand{\lcm}{\operatorname{lcm}}
\newcommand{\SNF}{\operatorname{SNF}}
\newcommand{\Jac}{\operatorname{Jac}}
\newcommand{\rank}{\operatorname{rank}}
\newcommand{\Ind}{\operatorname{Index}}
\newcommand{\disc}{\operatorname{disc}}
\begin{document}
\begin{center}
{\LARGE\bfseries The Gamma Identities}\\[0.4em]
{\large the reflection formula and the Gauss multiplication formula}\\[0.6em]
{\normalsize Isaev Iskhak Khamzatovich}\\[0.2em]
{\normalsize The <<Dynamic Principle>> Program --- the \texttt{hodge-laboratory}}\\[0.15em]
{\small Standalone theorem monograph · Монография 3/16}
\end{center}
\vspace{0.4em}
{\color{linkblue}\hrule height 0.8pt}
\vspace{0.8em}

\section{Setting}

The entire analytic toolkit of the $N=15/30$ stand consists of two
functional identities of the gamma function. Nothing else is needed:
all relations among the $\Ga$-products used by the program follow
from these two.

\begin{lemma}[reflection formula]\label{lem:reflection}
For all $z\in\CC\setminus\ZZ$,
$\Ga(z)\Ga(1-z)=\frac{\pi}{\sin\pi z}$.
\end{lemma}

\begin{proof}
The key is Euler's integral
$\int_0^\infty\frac{u^{z-1}}{1+u}du=\frac{\pi}{\sin\pi z}$
($0<\re z<1$), proved by residues over the keyhole contour: the
unique pole $u=-1$ has residue $1$; the two sides of the cut give
$\bigl(e^{i\pi(1-z)}-e^{-i\pi(1-z)}\bigr)\int_0^\infty
\frac{x^{z-1}}{1+x}dx=2\pi i$. The substitution $t=\frac{u}{1+u}$
turns the integral into $\Bfun(z,1-z)=\Ga(z)\Ga(1-z)$; analytic
continuation extends the identity.
\end{proof}

\begin{lemma}[Gauss multiplication]\label{lem:gauss}
For integers $m\ge2$ and $mz\notin\ZZ_{\le0}$,
$\prod_{k=0}^{m-1}\Ga(z+\tfrac km)=(2\pi)^{\frac{m-1}2}
m^{\frac12-mz}\Ga(mz)$.
\end{lemma}

\begin{proof}
The ratio $Q_m=G_m/H_m$ is $1$-periodic and holomorphic: the
recurrence $\Ga(w+1)=w\Ga(w)$ gives equal logarithmic derivatives.
Stirling's formula gives $Q_m\to1$ as $\re z\to+\infty$ in vertical
strips; periodicity spreads the limit; hence $Q_m\equiv1$.
\end{proof}

\begin{lemma}[inventory of relations]\label{lem:relations}
All used relations among $\Omega_{a,b}=\Bfun(a/N,b/N)$ follow from
the two identities: (1) $a+b=N$: $\Omega_{a,b}=\pi/\sin(\pi a/N)$;
(2) $a+b>N$: reduction
$\Omega_{a,b}=\frac{N}{a+b-N}\frac{\Ga(a/N)\Ga(b/N)}
{\Ga((a+b-N)/N)}$; (3) $a+b<N$: canonical.
\end{lemma}

\begin{proof}
Point 1 --- reflection at $b/N=1-a/N$; point 2 --- the recurrence
$\Ga(\frac{a+b}N)=\frac{a+b-N}N\Ga(\frac{a+b-N}N)$; point 3 ---
by definition.
\end{proof}

\begin{databox}[numbers]
The reflection ladder (V6): $\Ga(k/N)\Ga(1-k/N)=\pi/\sin(\pi k/N)$
for all $k$, $N=15,30$; maximal relative deviation
$7.0\cdot10^{-36}$ at $dps=35$.
\end{databox}

\begin{statusbox}[summary]
Both identities are derived from integral representations and the
residue theorem --- no external results. The inventory fixes the
scope: the full classification of $\QQ$-relations (Rohrlich) is not
used --- sufficiency is verified by the rank (V7: rank $=g$).
\end{statusbox}

\begin{verifybox}[reproduction]
\texttt{python3 laboratory.py} $\to$ ``Stand N=15/30 $\to$ V6'';
the plot \texttt{fig\_gamma.png} ($600$~dpi) is the numerical
portrait of the identity.
\end{verifybox}

\vfill
{\color{linkblue}\hrule height 0.6pt}
\vspace{0.5em}
{\small\color{gray}
License: individual exclusive license of Isaev Iskhak Khamzatovich.
All rights to the computations, formulas, and texts belong to the author.
Verification: \texttt{python3 laboratory.py} (``Stands''/``Certificates''),
Lean: \texttt{verification/lean/HodgeLaboratory.lean}.
}
\end{document}
